\clearpage
\SourceLeaf{356}
\pdfbookmark[2]{8. Gauss to Olbers, 29 May 1815}{parabolisch-008}
\begin{center}
[8.]\par
\medskip
\textsc{Gauss} to \textsc{Olbers}, Göttingen, 29 May 1815.
\end{center}
\noindent\hspace*{\parindent}.\ .\ .\ .\ .\ .\quad I am somewhat inclined to publish my theoretical investigations into the calculation of cometary orbits in general someday in a separate work, as a supplement to my Th. M. C. C. It might run to perhaps 6—8 sheets, and I would then add a table for parabolic motion arranged in a new way, somewhat easier to use than \textsc{Barker}'s. This might take up another 3 sheets. The groundwork has already been laid, and Mr. \textsc{Encke} has helped me with it.\quad .\ .\ .\ .\ .
\par\medskip
\begin{center}\rule{.18\textwidth}{.4pt}\end{center}
\medskip
\pdfbookmark[2]{9. Gauss to Bessel, 24 June 1815}{parabolisch-009}
\begin{center}
[9.]\par
\medskip
\textsc{Gauss} to \textsc{Bessel}, Göttingen, 24 June 1815.
\end{center}
\noindent\hspace*{\parindent}.\ .\ .\ .\ .\ .\quad I am half inclined to bring out a supplement to the Theoria motus corporum coelestium for parabolic orbits before long. I have several theoretical additions for it that are of some consequence. I also intend to include a new table for parabolic motion, arranged even more conveniently than \textsc{Barker}'s. Its calculation has already progressed so far that I no longer use any other table myself.\quad .\ .\ .\ .\ .
\par\medskip
\begin{center}\rule{.22\textwidth}{.4pt}\end{center}
